\section{Introduction}

% Plan:
% - introduce HAR and cite lara and reyes-ortiz (base papers for HAR)
% - add background information on sensors being faulty (different kinds of faults)
% - works that investigated missing sensor data cite saha and hossain (papers on missing data)
% - plan of what we investigated (corruption framework) and some of the models we used
% - answer question: what is the point of this paper?
%   - answer: provide a baseline for sensor corruption covering uncovered corruption types
%   - could have a stronger answer/conclusion but this will depend on the experiments and results

Human Activity Recognition (HAR) is a very active area of research due to technological
advancements of smartphones, microelectronics, and sensors. This is the purpose of
\textit{Ubiquitous sensing}, a research area focused of collecting data acquired from sensors and
utilizing it to extract information \cite{lara-2013}. The recognition of human activities has
many uses in healthcare, monitoring, and human-computer interaction \cite{lara-2013}.
Reyes-Ortiz et al. \cite{reyes-ortiz-2013} provided the Human Activity Recognition Using
Smartphones (UCI HAR) dataset, a benchmark built from accelerometer and gyroscope signals collected
from 30 subjects performing activities. Subsequent deep learning approaches have matched or
exceeded their results on clean data \cite{ronao-2016}.

Despite the strong performance, real-world sensor systems regularly produce degraded data. Sensors
age, drift, and miscalibrate, for example: wireless transmission introduces packet loss
\cite{ni-2009}; battery depletion or physical damage can cause intermittent failure \cite{ni-2009}.
The resulting corruption spans a broad range of failures that go beyond simple missing data,
including constant bias offsets, gradual temporal drift, amplification/attenuation gain shifts,
etc. Understanding how HAR models respond to these different forms of degradation is important for
building reliable activity recognition systems that generalize beyond clean data.

A number of works have investigated the problem of missing sensor data in HAR \cite{hossain-2018,
hossain-2020, saha-2025}, demonstrating that random dropout can substaintially degrade
classification accuracy. However these works focus on missing sensor data only, missing other
sensor fault types. To the best of our knowledge, no prior work has provided a systematic study of
how structured sensor corruption, varying by type, severity, and granularity (sensor modality)
affects HAR classifier performance.

In this work, we address this gap by introducing a structured corruption framework defined along
three dimensions: corruption type, granularity, and severity. Our framework applies five sensor
fault types: stochastic noise, dropout, bias, gain (amplification, and attenuation), and drift, at
controlled severity levels and at granularities ranging from individual signal axes, to full sensor
modalities. We apply this framework to the UCI HAR dataset, training four classifiers on clean and
corrupted data, and evaluating each under systematic corruption. We thus provide the first baseline
for understanding how different categories of sensor corruption affect activity recognition
performance.
